% Pregunta 2. Mercado de trabajo y productividad (bloque II.1) · peso alto
\clearpage
\section*{Pregunta 2. Mercado de trabajo y productividad}

\begin{enunciado}
\textbf{2.} Caracterice el \textbf{mercado laboral} español a lo largo de la serie. ¿Puede afirmarse que se encuentra tensionado? Analice la evolución de la curva de Beveridge. ¿A qué podría deberse la caída de la temporalidad desde 2022? Valore la evolución de la \textbf{productividad} y su contribución al crecimiento.

\smallskip
\textbf{Peso alto} · \tiempo{15} · guía: II.1 (cantidades, productividad, CLU y margen, tipos de paro, \autor{Okun}, \autor{Beveridge}), I.5 y V.5 · cuadros 3 y 2~bis; gráfico 3.
\end{enunciado}

\begin{examen}
\textit{El mercado de trabajo español amplifica el ciclo: destruye empleo con gran intensidad en las recesiones y lo crea con la misma rapidez en las expansiones. Hoy está en máximos de empleo y de actividad, con el paro más bajo desde 2008, pero conserva sus 2 debilidades de fondo: un paro estructural que dobla al de la zona del euro y una productividad que apenas crece.}

\rotulo{a) Caracterización a lo largo de la serie}
La tasa de paro (cuadro~3) dibuja el ciclo con nitidez: del 8,2~\% en 2007 al 26,1~\% en 2013, para bajar al 14,1~\% en 2019 y al 10,5~\% en 2025 (10,2~\% en el primer semestre de 2026, gráfico~3). Es todavía el doble que en la zona del euro (6,4~\%), con un paro juvenil del 24,9~\% (55,5~\% en 2013). La tasa de actividad está en máximos (74,9~\%) y la de empleo (67,0~\%) supera ya la de 2007: la mejora no se debe a un \emph{efecto desánimo}, sino a la creación de empleo.

Lo que más llama la atención es \textbf{cómo se ajusta}. En la crisis el ajuste fue por \emph{cantidades}: en 2009 el empleo cayó un 6,7~\% con un PIB que caía un 3,8~\%, mientras la remuneración por asalariado aún crecía un 4,0~\%; los salarios solo se moderaron después (en torno al 0~\% en 2011-2014), y los costes laborales unitarios cayeron en 2010-2014, una \emph{devaluación interna} tardía. Es la huella de un mercado dual, en el que los temporales absorbían el choque. En 2020, en cambio, el empleo cayó un 2,9~\% y las horas un 11,0~\%: el ajuste se hizo con las horas, lo que sería coherente con esquemas de mantenimiento del empleo que, en lugar de despedir, reducían la jornada.

La relación con el ciclo es intensa. En 2014-2019 el paro bajó 2,0 puntos al año con un crecimiento medio del 2,6~\%; en 2022-2025 solo 1,1 puntos con un 3,8~\%. Ahora hace falta más crecimiento para bajar el paro porque la población activa crece con la inmigración; medido frente al potencial, la respuesta del paro es parecida en las 2 fases (en torno a 1 punto de paro por punto de crecimiento por encima del potencial, el doble que el medio punto de la ley de \autor{Okun} en EE.~UU.). Y en 2023-2025 el paro bajó casi un punto al año creciendo al ritmo del potencial o algo por debajo, lo que apuntaría a una reducción del paro estructural. Ayuda que el paro de larga duración haya caído del 13,0~\% de los activos en 2013 al 3,4~\%: la crisis dejó menos \emph{histéresis} de la que cabía temer, aunque el paro, que en 2019 aún estaba en el 14~\%, sigue muy por encima de la tasa de pleno empleo de una economía avanzada.

\rotulo{b) ¿Está tensionado?}
\textbf{Algo más que antes, pero no en conjunto.} A favor de la tensión: la tasa de vacantes está en su máximo desde 2011 (0,9~\%), la relación entre vacantes y parados ha subido de 0,05 en 2019 a casi 0,09 en 2025, y la remuneración por asalariado crece un 4,4~\% con una productividad estancada, de modo que los costes laborales unitarios suben el 4,4~\% (3,3~\% en la zona del euro). En contra: una economía que produce en torno a su potencial (brecha del 0,6~\%), un paro del 10,5~\%, un paro juvenil del 25~\% y unas vacantes que siguen siendo muy pocas frente a las de la zona del euro. La tensión parece, por tanto, \emph{localizada} (sectores, ocupaciones y territorios donde falta mano de obra) más que general; el riesgo es que alimente los efectos de segunda ronda del choque energético (pregunta~4).

\rotulo{c) La curva de \autor{Beveridge}}
El gráfico~3 recorre 3 tramos. Entre 2010 y 2013 la economía baja a lo largo de la curva hacia la derecha (vacantes del 1,25 al 0,57~\%, paro del 19,9 al 26,1~\%): es el ciclo, la recesión. Entre 2014 y 2019 deshace el camino hacia arriba y a la izquierda (el paro baja al 14,1~\% y las vacantes suben al 0,70~\%): la recuperación. Tras el paréntesis de 2020, desde 2021 sigue hacia la izquierda, y desde 2022 lo hace casi en horizontal: \uline{el paro baja del 13,0 al 10,2~\% con las vacantes estables en el 0,9~\%}. Eso sería compatible con un desplazamiento de la curva hacia el origen, es decir, con un mejor emparejamiento entre parados y puestos, quizá favorecido por el cambio normativo que se comenta en el apartado~d) y por la llegada de trabajadores a los sectores con escasez. En todo caso, la curva española sigue muy alejada de la europea: muy pocas vacantes con mucho paro, la combinación de un paro estructural alto y de problemas de emparejamiento (cualificaciones, territorio, vivienda); conviene leerla por su evolución más que por su nivel, porque la medición de las vacantes también influye.

\rotulo{d) La caída de la temporalidad desde 2022}
Los asalariados temporales pasan del 25,4~\% en 2021 al 15,4~\% en 2025: 10 puntos en 4 años y \emph{en plena expansión}, cuando lo habitual en España era lo contrario (la temporalidad subió del 24,0 al 26,9~\% entre 2014 y 2018). Una caída tan brusca, persistente y a contracorriente del ciclo no parece cíclica: apunta a un \textbf{cambio normativo} que haya restringido la contratación temporal a finales de 2021. Con matices: parte del trasvase puede ser a contratos indefinidos menos estables (discontinuos o de corta duración), que la tasa de temporalidad no recoge como empleo inestable. Si se consolida, reduce la dualidad, da más estabilidad al consumo y más incentivos a formar a los trabajadores, aunque su efecto sobre la productividad está por ver.

\rotulo{e) La productividad y su contribución al crecimiento}
La productividad es \textbf{el punto débil}. Por hora trabajada crece solo un 0,33~\% anual entre 2019 y 2025 (1,9~\% acumulado) y por ocupado cae ligeramente ($-0{,}13~\%$ anual), porque se trabajan menos horas por persona (cuadros~2~bis y~3). De los 1,59~puntos de crecimiento medio del PIB, la productividad por hora aporta 0,33, una quinta parte: el resto lo pone el trabajo. En 2007-2013 la productividad creció más (1,6~\% anual), pero de forma espuria, porque se destruyó sobre todo el empleo menos productivo (construcción, temporales): en España la productividad aparente suele ser \emph{contracíclica}.

Como $\gamma_{\mathrm{CLU}}=\gamma_W-\gamma_{\mathrm{PAT}}$, con una productividad estancada todo el aumento salarial pasa a los costes: en 2025, $4{,}4-0{,}0=4{,}4~\%$. Las causas de fondo serían el peso de los sectores de baja productividad, el pequeño tamaño medio de las empresas, el escaso esfuerzo en I+D (1,50~\% del PIB, frente al 2,26~\% de la UE, cuadro~6) y una inversión en maquinaria que apenas crece desde 2019. Con una población que envejece (la tasa de dependencia de los mayores pasa del 24,0 al 31,2~\%), \uline{el crecimiento futuro dependerá cada vez más de la productividad y no de sumar trabajadores}: es el reto que justifica las reformas de la pregunta~7.
\end{examen}

\begin{complemento}{2.1 · Por qué se contesta así: el tipo de paro decide la receta}
\que{Qué se aplica.} Guía, II.1: los rótulos con cada subpregunta al principio (nota al opositor de II.1: «para que el tribunal las encuentre al oírlas»); la tasa de actividad para descartar el efecto desánimo; la flexibilidad del mercado (¿se ajusta por precios o por cantidades?); el tipo de paro (gráfico~\ref{fig:tipos-paro}); \autor{Okun} con $\Delta U$; \autor{Beveridge}; y la cadena productividad-CLU-margen.

\que{Un dato de fuera que encaja.} La Comisión estima una NAWRU de España en torno al 12~\% \parencite{CE2026AMECO} (guía, II.1). Con un paro del 10,5~\%, la economía estaría por debajo de su paro de equilibrio, lo que cuadraría con unos CLU que se aceleran; el FMI, en cambio, sitúa el paro estructural en torno al 10~\%, es decir, que el paro actual ya estaría en su nivel estructural \fmi{6}, coherente con la brecha casi nula del cuadernillo. Es la respuesta de máximos a «¿está tensionado?»: \emph{para su propio paro estructural, sí; para el de la zona del euro, no}.

\begin{figuracaja}
\caption{Tipos de paro: cómo se detectan y qué receta tienen (infografía de la guía)}\label{fig:tipos-paro}
\grafico[0.52\linewidth]{esquema-tipos-paro}
\fuente{Guía del dictamen, II.1. Qué hay que ver: cada tipo deja una huella distinta en el cuadro; la tabla siguiente la busca en el cuadernillo de España.}
\end{figuracaja}

\begin{tablacaja}
\centering\scriptsize
\caption{¿Qué paro tiene España? Los indicios del cuadernillo}\label{tab:tipo-paro}
\begin{tabular}{>{\raggedright\arraybackslash}p{2.6cm}>{\raggedright\arraybackslash}p{6.2cm}>{\raggedright\arraybackslash}p{5.2cm}}
\toprule
\rowcolor{tceeazulpalido}
\textbf{Tipo de paro} & \textbf{Indicio en el cuadernillo} & \textbf{Lectura} \\
\midrule
Keynesiano (demanda) & Consumo privado +3,1~\% y brecha +0,6~\% en 2025 & \textcolor{tceeverde}{No}: la demanda tira \\
Clásico (salarios) & CLU +4,4~\% (zona del euro 3,3~\%); CLU reales +1,4 en 2025; sin embargo, de 2007 a 2025 los CLU crecen menos que en la zona & \textcolor{tceedorado}{Riesgo reciente}, no un problema acumulado \\
Friccional o de emparejamiento & Vacantes muy bajas con paro alto; paro juvenil del 24,9~\%; Beveridge lejos de la europea & \textcolor{tceerojo}{Sí}: el problema principal \\
De larga duración & 3,4~\% de los activos (13,0~\% en 2013) & \textcolor{tceeverde}{Mejora mucho}; histéresis menor \\
\bottomrule
\end{tabular}
\fuente{Cuadros 1, 3 y 6 y gráfico 3 del cuadernillo. Recetas: guía, V.5 (políticas activas, formación, movilidad y vivienda para el emparejamiento).}
\end{tablacaja}

\end{complemento}

\begin{complemento}{2.2 · La ley de \autor{Okun} con los datos}
\que{La fórmula de la guía} (II.1): $\Delta U_t\approx-0{,}5\,(\gamma_{Y,t}-\gamma_{Y^p,t})$, con $U$ la tasa de paro. Se aplica con el potencial del cuadro~1, no con el crecimiento a secas.


\que{Cálculo}: 2014-2019, $\Delta U$ medio $=-2{,}0$; $\gamma_Y-\gamma_{Y^p}=2{,}62-0{,}96=1{,}7$; coeficiente $-2{,}0/1{,}7\approx-1{,}2$. 2022-2025, $\Delta U=-1{,}1$; $\gamma_Y-\gamma_{Y^p}=3{,}75-2{,}59=1{,}2$; coeficiente $\approx-0{,}95$. \uline{El coeficiente casi no cambia: lo que cambia es el potencial}, que ahora incluye la llegada de activos.

\que{Por qué baja el paro con crecimiento igual al potencial.} Si el potencial del FMI incorpora una NAWRU que baja, el paro observado baja con ella aunque la brecha no cambie. Es la lectura optimista (reforma laboral, mejor emparejamiento); la prudente, que el potencial esté sobrestimado. Las 2 se dicen con condicional.

\que{Error típico.} Decir que «la ley de \autor{Okun} no se cumple» porque en 2022-2025 el paro baja menos con más crecimiento: la guía advierte que el coeficiente varía y que hay que comparar con el potencial (II.1).

\begin{figuracaja}
\caption{España: variación del paro y crecimiento por encima del potencial, 2008-2025}\label{fig:s03}
\grafico[0.74\linewidth]{s03-okun}
\fuente{Cuadros~1 y~3 del cuadernillo (FMI, WEO de octubre de 2026: PIB, potencial y paro). Qué hay que ver: los puntos de 2014-2019 caen por debajo de la recta de \autor{Okun} de EE.~UU.: en España el paro reacciona más al crecimiento; en 2023-2025 (crecimiento igual o algo inferior al potencial) el paro baja casi un punto al año; 2020-2021 se salen de la relación (ERTE).}
\end{figuracaja}
\end{complemento}

\begin{complemento}{2.3 · \autor{Beveridge}, temporalidad y productividad: cálculos y contexto}
\que{Tensión: vacantes por parado} (guía, II.1: «la ratio vacantes/parados resume la tensión»): $0{,}70/14{,}1=0{,}050$ en 2019; $0{,}90/10{,}5=0{,}086$ en 2025; $0{,}90/10{,}25=0{,}088$ en el primer semestre de 2026. En la zona del euro, $2{,}25/6{,}3=0{,}36$ en 2025, 4 veces más \parencite{Eurostat2026jvs} (dato de la guía, no del cuadernillo).

\que{La serie de vacantes tiene un salto} entre 2009 y 2010 (el cuadernillo la da desde 2010): no se compara 2009 con 2010.

\que{Temporalidad y ciclo.} Antes de 2022 era procíclica: bajó en la crisis (del 31,6 al 23,2~\% entre 2007 y 2013, porque se destruía sobre todo empleo temporal) y subió en la recuperación. Su caída en 2022-2025, con el PIB creciendo, rompe el patrón: \emph{un cambio de patrón que coincide con un cambio normativo} es la prueba indirecta que el tribunal espera.

\que{Productividad, en 3 periodos} (cuadro~3, acumulada): por hora, +10,2~\% en 2007-2013, +2,7~\% en 2013-2019 y +1,9~\% en 2019-2025; por ocupado, +9,8, +2,2 y $-0{,}8~\%$. \uline{La que más crece es la de la crisis}: efecto composición (guía, I.5, «cuidado»).

\que{Salario real y productividad, 2019-2025}: la remuneración por asalariado sube un 25,8~\% y el IPCA un 21,5~\%: salario real de consumo +3,5~\%, por encima de la productividad por ocupado ($-0{,}8~\%$) y por hora (+1,9~\%). Es otra forma de ver por qué los CLU reales suben en 2025 y los márgenes se estrechan (pregunta~4).

\que{Alternativas de redacción.} Para la tensión: «el mercado está más tensionado que en 2019, pero lejos de la tensión europea: la escasez de mano de obra convive con un paro alto, la marca de un problema de emparejamiento». Para la temporalidad, si no se quiere nombrar la reforma: «un cambio institucional que limita la contratación temporal».

\que{Qué añadir con más tiempo.} La demografía (guía, B.1): el envejecimiento resta crecimiento potencial y la inmigración lo suma, con efectos sobre el paro a corto plazo (la inmigración eleva primero el paro y después el potencial). Y el salario mínimo, que el cuadernillo no trae: si se menciona, sin cifras y con condicional.

\que{Qué recortar en el examen.} La frase de diagnóstico; el paro en 3 cifras (2007, 2013, 2025) frente a la zona; el ajuste por cantidades en 1 frase; tensión en 2 frases; \autor{Beveridge} en 3 tramos; temporalidad en 2 frases; productividad con la cifra de 0,33 y el CLU de 2025.
\end{complemento}

\begin{complemento}{2.4 · Contexto y fuentes: lo que el cuadernillo no trae}
\que{La reforma laboral.} El cambio normativo que el texto de examen deduce de los datos es el Real Decreto-ley 32/2021, de 28 de diciembre, de medidas urgentes para la reforma laboral: limitó los contratos temporales (causas tasadas, duraciones máximas) y reforzó el contrato fijo discontinuo. Por eso la tasa de temporalidad cae en 2022, el primer año completo de aplicación. Nombrarla en el examen es razonable porque el país está identificado, pero la prueba que puntúa es la de los datos: \emph{una caída brusca, persistente y a contracorriente del ciclo}. El debate abierto (que se dice con condicional) es cuánto empleo pasó a fijos discontinuos, que la EPA cuenta como indefinidos.

\que{Los esquemas de mantenimiento del empleo.} Los expedientes de regulación temporal de empleo (ERTE) de la pandemia (Real Decreto-ley 8/2020, de 17 de marzo) permitieron suspender contratos o reducir jornada con prestación pública: explican que en 2020 el empleo cayera un 2,9~\% y las horas un 11,0~\% (cuadro~3), y que la guía (II.1) advierta que en la pandemia el paro apenas reflejó la caída del PIB.

\que{La NAWRU y la curva europea.} La NAWRU de la Comisión para España (en torno al 12~\%, \cite{CE2026AMECO}) y la curva de \autor{Beveridge} de la zona del euro (vacantes del 2,25~\% con un paro del 6,3~\% en 2025, \cite{Eurostat2026jvs}) están en los cuadros 2.1 y 2.3; no en el texto de examen, porque no vienen en el cuadernillo.

\que{Las vacantes, con cautela.} La serie de vacantes de Eurostat para España salta entre 2009 y 2010 sin marca de ruptura (por eso el cuadernillo la da desde 2010): se leen sus cambios, no su nivel frente a otros países.
\end{complemento}

\begin{complemento}{2.5 · El diagnóstico del FMI (artículo~IV de 2026) frente a los datos}
\que{Paro estructural.} El paro «parece estabilizarse en un nivel estructural de en torno al 10~\%» \fmi{6}; el paro estructural habría bajado en la última década, como sugieren el desplazamiento de la curva de \autor{Beveridge} hacia la izquierda y la ausencia de aceleración de los salarios reales pese a la caída continua del paro \fmi{27}. \textbf{Coincide} con la lectura de la solución (curva que se desplaza hacia el origen, paro estructural que baja) y aporta una cifra que el cuadernillo no da. Con un paro del 10,5~\%, la economía estaría casi en su paro estructural, coherente con la brecha casi nula.

\que{Salarios.} Los salarios pactados en convenio crecían en torno al 3,5~\% a finales de 2025 y los efectivos por asalariado, un 4,9~\%: la diferencia (\emph{deriva salarial}) refleja un mercado tensionado y el intento de recuperar la inflación pasada \fmi{7}. Es la base del riesgo de segunda ronda de la pregunta~4.

\que{La reforma de 2021.} Según el FMI, redujo de forma duradera la temporalidad y mejoró la estabilidad del empleo, aunque «modestamente», porque bajó la duración de los contratos indefinidos y subió el peso de los fijos discontinuos \fmi{27}: \textbf{coincide} con el matiz del texto de examen.

\que{Recomendaciones del FMI} \fmi{28-29}: cotizaciones más altas para las empresas que despiden y contratan con más frecuencia (\emph{experience rating}); financiar los servicios públicos de empleo según su tasa de colocación; ampliar más allá de 6 meses la compatibilidad del subsidio con el empleo; no subir el salario mínimo por encima del objetivo del 60~\% del salario medio neto ni indexarlo al IPC; y un crédito fiscal por trabajo para los hogares de renta baja. El FMI subraya además que la tasa de empleo sigue siendo de las más bajas de la zona del euro \fmicp{1-2}.

\que{Dónde difiere.} El FMI duda de que el paro pueda bajar mucho más sin nuevas reformas que no están sobre la mesa \fmi{27}; la solución es algo más optimista al leer la caída del paro de 2023-2025 con un crecimiento igual al potencial como reducción del paro estructural. Las 2 lecturas son compatibles: el paro estructural bajó, pero hasta el 10~\%.
\end{complemento}
